\exercisesection{}

除非另有说明，以下巴拿赫空间与希尔伯特空间均假定为复空间。

\begin{enumerate}
\def\labelenumi{\arabic{enumi})}
\tightlist
\item
  \begin{enumerate}
  \def\labelenumii{\alph{enumii})}
  \tightlist
  \item
    存在向量空间 E、F、G，从 E 到 F 的部分算子 u、v，以及从 F 到 G 的部分算子 w，使得 \(w \circ (u + v) \neq w \circ u + w \circ v\) 。
  \end{enumerate}
\end{enumerate}

\begin{enumerate}
\def\labelenumi{\alph{enumi})}
\setcounter{enumi}{1}
\item
  给出拓扑向量空间上一个可闭算子 u 的例子，使得 u 的闭包的定义域严格包含于 u 的定义域的闭包之中。
\item
  给出希尔伯特空间上闭算子 u 与 v 的一个例子，使得 u 连续但 \(u \circ v\) 不是闭的。
\end{enumerate}

\begin{enumerate}
\def\labelenumi{\arabic{enumi})}
\setcounter{enumi}{1}
\tightlist
\item
  设 E 与 F 为巴拿赫空间，G 为 E 的一个子空间。设 u 为从 G 到 F 的有限秩线性映射。假设 u 在 G 上不连续。
\end{enumerate}

\begin{enumerate}
\def\labelenumi{\alph{enumi})}
\item
  由 u 定义的、以 G 为定义域的从 E 到 F 的部分算子是不可闭的。
\item
  假设 E 与 G 为希尔伯特空间。计算由 u 定义的部分算子的伴随。
\end{enumerate}

\begin{enumerate}
\def\labelenumi{\arabic{enumi})}
\setcounter{enumi}{2}
\tightlist
\item
  设 K = R 或 C。设 E、F 为 K 上的拓扑向量空间，u 为从 E 到 F 的一个稠定算子。设 D' 为满足下列条件的连续线性形式 \(\lambda \in F'\) 组成的集合：存在连续线性形式 \(\mu \in E'\) 在 u 的定义域上与 \(\lambda \circ u\) 重合。
\end{enumerate}

\begin{enumerate}
\def\labelenumi{\alph{enumi})}
\tightlist
\item
  集合 D' 是 F' 的向量子空间，线性形式 \(\mu\) 由 \(\lambda\) 唯一确定，且从 D' 到 E' 的映射 \(\lambda \mapsto \mu\) 是线性的。
\end{enumerate}

我们把 u 的转置记作 \(^{t}u\) ，它是从 F' 到 E' 、以 D' 为定义域、由 \(\lambda\mapsto\mu\)（\(\lambda\in D'\) ）给出的部分算子。

\begin{enumerate}
\def\labelenumi{\alph{enumi})}
\setcounter{enumi}{1}
\item
  假设 E 与 F 是局部凸的。算子 u 可闭当且仅当对偶 \(\mathrm{F} \times \mathrm{dom}(^{t}u) \to \mathrm{K}\) 是分离的。
\item
  假设 E 与 F 是局部凸的且 F 是半自反的。算子 u 可闭当且仅当 \(^{t}u\) 稠定。
\end{enumerate}

\begin{enumerate}
\def\labelenumi{\arabic{enumi})}
\setcounter{enumi}{3}
\item
  设 E 为复希尔伯特空间，u 为 E 上的一个单射自伴部分算子。部分算子 \(u^{-1}\) 是自伴的。
\item
  设 E 为复巴拿赫空间。E 上的连续半群是指一个映射 \(s\colon\mathbf{R}_{+}\to\mathcal{L}(\mathrm{E})\) ，它满足
\end{enumerate}

\begin{enumerate}
\def\labelenumi{(\roman{enumi})}
\item
  \(s(0) = 1_{\mathrm{E}}\) ；
\item
  从 \(R_{+} \times E \to E\) 、由 \((t, x) \mapsto s(t)x\) 定义的映射是连续的；
\item
  对 \(t_1\) 与 \(t_2\) （均属于 \(\mathbf{R}_{+}\) ） ，有 \(s(t_1 + t_2) = s(t_1)s(t_2)\) 。
\end{enumerate}

我们考虑 E 上的一个连续半群 s。

\begin{enumerate}
\def\labelenumi{\alph{enumi})}
\tightlist
\item
  存在实数 C ≥ 0 与 M ≥ 0，使得
\end{enumerate}

\[
\| s (t) \| \leqslant \mathrm{Ce} ^ {\mathrm{Mt}}
\]

对所有 \(t \in R_{+}\) 成立。（设 a \textgreater{} 0；证明 \(\|s(t)\|\) 在 \(t \in [0, a]\) 上有界。）在本习题的其余部分，我们固定这样的数。

\begin{enumerate}
\def\labelenumi{\alph{enumi})}
\setcounter{enumi}{1}
\item
  设 D 为 \(x \in E\) 中使得映射 \(\psi_{x}\) （从 \(R_{+}\) 到 E 、由 \(t \mapsto s(t)x\) 定义）在 0 处可微的那些 x 组成的集合。空间 D 在 E 中稠密；由 \(u(x) = -\psi_{x}'(0)\) 定义的从 D 到 E 的映射 u 是 E 上的一个闭部分算子，称为半群 s 的无穷小生成元。
\item
  设 \(\lambda\) 为使 \(\mathcal{R}(\lambda) < -\mathrm{M}\) 的复数。公式
\end{enumerate}

\[
v _ {\lambda} = - \int_ {\mathbf {R} _ {+}} e ^ {\lambda t} s (t) d t
\]

定义 E 的一个自同态。

\(d)\) \(v_{\lambda}\) 的像包含于 \(\operatorname{dom}(u)\) ，并且我们有 \(u \circ v_{\lambda} = -1_{\mathrm{E}} - \lambda v_{\lambda}\) 。

\begin{enumerate}
\def\labelenumi{\alph{enumi})}
\setcounter{enumi}{4}
\tightlist
\item
  我们有 \(u \circ v_{\lambda} = v_{\lambda} \circ u\) 。
\end{enumerate}

\(f)\) 复数 \(\lambda\) 属于 \(u\) 的预解集，并且我们有 \(v_{\lambda} = -\mathrm{R}(u,\lambda)\) 。

\begin{enumerate}
\def\labelenumi{\alph{enumi})}
\setcounter{enumi}{6}
\tightlist
\item
  我们有
\end{enumerate}

\[
\| v _ {\lambda} ^ {n} \| \leqslant \frac {\mathrm{C}}{| \mathcal {R} (\lambda) + \mathrm{M} | ^ {n}}
\]

对所有 \(n \in N\) 成立。

\begin{enumerate}
\def\labelenumi{\arabic{enumi})}
\setcounter{enumi}{5}
\tightlist
\item
  设 E 为复巴拿赫空间，u 为 E 上的一个稠定闭部分算子。假设存在实数 \(C \geqslant 0\) 与 \(M \geqslant 0\) ，使得
\end{enumerate}

\begin{enumerate}
\def\labelenumi{(\roman{enumi})}
\item
  集合 {]}−∞, −M{[} 包含于 u 的预解集；
\item
  对每个整数 \(n \in \mathbf{N}\) 与每个实数 \(\lambda < -\mathrm{M}\) ，我们有
\end{enumerate}

\[
\left\| \mathrm{R} (u, \lambda) ^ {n} \right\| \leqslant \mathrm{C} | \lambda + \mathrm{M} | ^ {- n}.
\]

\begin{enumerate}
\def\labelenumi{\alph{enumi})}
\tightlist
\item
  对每个实数 \(y < -\mathrm{M}\) ，我们记
\end{enumerate}

\[
u _ {y} = - y 1 _ {\mathrm{E}} + y ^ {2} \mathrm{R} (u, y) \in \mathcal {L} (\mathrm{E}).
\]

诸自同态 \(u_{y}\) 两两可交换；对每个 \(x \in \text{dom}(u)\) ，我们有 \(u_{y}(x) \to u(x)\) （\(y \to -\infty\) ）。

\begin{enumerate}
\def\labelenumi{\alph{enumi})}
\setcounter{enumi}{1}
\item
  从 \(R_{+}\) 到 \(\mathcal{L}(E)\) 、由 \(s_{y}(t)=e^{-tu_{y}}\) 定义的映射是 E 上的连续半群，使得 \(\|s_{y}(t)\|\leqslant Ce^{Mt}\) 对所有 \(t\in R_{+}\) 成立。
\item
  对所有 \(y_{1}, y_{2}\) 与 \(t\) （属于 \(\mathbf{R}_{+}^{*}\) ），以及每个 \(x\) （属于 \(\mathrm{dom}(u)\) ），我们有
\end{enumerate}

\[
\left\| s _ {y _ {1}} (t) x - s _ {y _ {2}} (t) x \right\| \leqslant t \mathrm{C} ^ {2} \left\| u _ {y _ {1}} (x) - u _ {y _ {2}} (x) \right\|.
\]

\(d)\) 设 \(t \in \mathbf{R}_{+}\) 。对 E 中每个 \(x\) ，定义在 \(\mathbf{R}_{+}^{*}\) 上的函数 \(y \mapsto s_{y}(t)x\) 当 \(y \to -\infty\) 时收敛于 E 的一个元素 \(s(t)x\) 。映射 \(s(t)\) 线性且连续，范数 \(\leqslant \mathrm{Ce}^{\mathrm{M}t}\) 。

\begin{enumerate}
\def\labelenumi{\alph{enumi})}
\setcounter{enumi}{4}
\item
  映射 \(t \mapsto s(t)\) 是 E 上的连续半群。
\item
  部分算子 u 是连续半群 s 的无穷小生成元（“希尔–吉田定理”）。（首先验证 s 的无穷小生成元是 u 的一个扩张。）
\end{enumerate}

\(g)\) 设 \(\widetilde{s}\) 为以 \(u\) 为无穷小生成元的一个连续半群。我们有 \(\widetilde{s} = s\) 。

\begin{enumerate}
\def\labelenumi{\arabic{enumi})}
\setcounter{enumi}{6}
\tightlist
\item
  设 E 为复巴拿赫空间，u 为 E 上的一个部分算子。我们称 u 是增生的，如果对每个 \(x \in \text{dom}(u)\) 存在连续线性形式 \(\ell \in E'\) 使得
\end{enumerate}

\begin{enumerate}
\def\labelenumi{(\roman{enumi})}
\item
  \(\mathcal{R}(\ell(u(x))) \geqslant 0;\)
\item
  我们有 \(\|\ell\|=\|x\|\) 且 \(\ell(x)=\|x\|^{2}\) 。
\end{enumerate}

我们称 E 上的连续半群 \(s\) 是压缩的，如果 \(\|s(t)\|\leqslant1\) 对所有 \(t\in\mathbf{R}_{+}\) 成立。

\begin{enumerate}
\def\labelenumi{\alph{enumi})}
\item
  部分算子 \(u\) 是 E 上一个连续压缩半群的无穷小生成元，当且仅当 \(u\) 是增生的并且存在 \(\lambda > 0\) 使部分算子 \(\lambda 1_{\mathrm{E}} + u\) 是满射的。（为证明条件必要，设 \(s\) 为以 \(u\) 为无穷小生成元的压缩半群；对 \(x \in \operatorname{dom}(u)\) ，考虑由 \(t \mapsto \ell(s(t)x)\) 定义的函数；它在 0 处可微，且其在 0 处的导数为 \(-\mathcal{R}(\ell(u(x)))\) 。）
\item
  假设 u 是闭的，且 \(u\) 与 \(^{t}u\) 都是增生的。则 u 是 E 上一个连续压缩半群的无穷小生成元。
\end{enumerate}

\begin{enumerate}
\def\labelenumi{\arabic{enumi})}
\setcounter{enumi}{7}
\tightlist
\item
  设 E 为复巴拿赫空间。设 s 为 E 上的一个连续压缩半群，u 为其无穷小生成元。我们有
\end{enumerate}

\[
\lim _ {n \to + \infty} \left(1 + \frac {t}{n} u\right) ^ {- n} x = s (t) x
\]

对所有 \(x \in E\) 成立。

\begin{enumerate}
\def\labelenumi{\arabic{enumi})}
\setcounter{enumi}{8}
\tightlist
\item
  \begin{enumerate}
  \def\labelenumii{\alph{enumii})}
  \tightlist
  \item
    求希尔伯特空间 E 与 F，以及从 E 到 F 的部分算子 u 与 v，使得 u 与 \(u + v\) 都稠定，且 \(u^{*} + v^{*}\) 不等于 \((u + v)^{*}\) 。
  \end{enumerate}
\end{enumerate}

\begin{enumerate}
\def\labelenumi{\alph{enumi})}
\setcounter{enumi}{1}
\item
  求希尔伯特空间 E、F、G，以及分别从 E 到 F 与从 F 到 G 的稠定部分算子 u 与 v，使得 \(v \circ u\) 稠定且 \(u^{*} \circ v^{*}\) 不等于 \((v \circ u)^{*}\) 。
\item
  求一个希尔伯特空间 E 与 E 上的自伴部分算子 u 与 v，使得 \(u + v\) 稠定但不是自伴的。
\end{enumerate}

\begin{enumerate}
\def\labelenumi{\arabic{enumi})}
\setcounter{enumi}{9}
\item
  给出乘法算子 \(m_{g_{1}}\) 与 \(m_{g_{2}}\) 的一个例子，使得 \(m_{g_{1}} \circ m_{g_{2}} \neq m_{g_{1}g_{2}}\) 。
\item
  我们把空间 \(\mathcal{D}(]0,1[)\) 与 \(\mathrm{L}^2 ([0,1])\) 的一个子空间等同起来。设 \(u\) 为 \(\mathrm{L}^2 ([0,1])\) 上以 \(\mathcal{D}(]0,1[)\) 为定义域、满足 \(u(f) = i^{-1}f'\) （对每个函数 \(f \in \mathcal{D}(]0,1[)\) ）的部分算子。
\end{enumerate}

\begin{enumerate}
\def\labelenumi{\alph{enumi})}
\item
  部分算子 \(u\) 是可闭的且 \(\overline{u}\) 是对称的。
\item
  \(\overline{u}\) 的属于 i 与属于 -i 的特征空间不是正交的。
\end{enumerate}

\begin{enumerate}
\def\labelenumi{\arabic{enumi})}
\setcounter{enumi}{11}
\item
  设 \(S \subset \mathbf{C}\) 为闭子集。给出一个稠定闭算子 \(u\) （定义在希尔伯特空间 \(E\) 上）的例子，使得 \(\operatorname{Sp}(u) = S\) 。
\item
  设 \(u\) 为希尔伯特空间 E 上的一个不可闭部分算子。不存在复数 \(\lambda\) 使得从 \(\text{dom}(u)\) 到 E 的、由 \(x \mapsto \lambda x - u(x)\) 定义的线性映射是双射且具有连续的逆。
\item
  设 \(u\) 为复希尔伯特空间 E 上的一个对称部分算子。映射 \(\lambda \mapsto \dim \operatorname{Ker}(\lambda 1_{\mathrm{E}} - u^{*})\) （从 \(\mathbf{C}\) 到 \(\overline{\mathbf{N}}\) ）在 \(u\) 的预解集上是局部常数的。
\item
  设 E 为复希尔伯特空间 \(\mathbf{C}^2\) ，\(u \in \mathcal{L}(\mathrm{E})\) 为满足 \(u(ae_1 + be_2) = be_1\) 的自同态。
\end{enumerate}

\begin{enumerate}
\def\labelenumi{\alph{enumi})}
\tightlist
\item
  证明 \(u\) 的谱退缩为 0，并且对所有 \(\lambda \in \mathbf{C}^*\) ，我们有公式
\end{enumerate}

\[
\| \mathrm{R} (u, \lambda) \| ^ {2} = \frac {2}{1 + 2 | \lambda | ^ {2} - \sqrt {1 + 4 | \lambda | ^ {2}}}.
\]

由此导出伪谱 \(\mathrm{PSp}_{\varepsilon}(u)\) （对所有 \(\varepsilon > 0\) ）的形状。

\begin{enumerate}
\def\labelenumi{\arabic{enumi})}
\setcounter{enumi}{15}
\item
  给出一个例子，说明 IV 第 251 页命题 21 的结论对谱的有界分支一般不成立。
\item
  \begin{enumerate}
  \def\labelenumii{\alph{enumii})}
  \tightlist
  \item
    设 U 为 C 的连通开子集，V ⊂ U 为开集。设 E 为复巴拿赫空间。设 \(f\colon\) U → E 为全纯映射，\(z_0 \in\) U 且 M \textgreater{} 0 满足：\(\|f(z)\| \leqslant\) M 对所有 \(z \in\) V 成立，且 \(\|f(z_0)\| <\) M 。则 \(\|f(z)\| <\) M 对所有 \(z \in\) U 成立。
  \end{enumerate}
\end{enumerate}

\begin{enumerate}
\def\labelenumi{\alph{enumi})}
\setcounter{enumi}{1}
\tightlist
\item
  设 E 为复巴拿赫空间，\(u \in \mathcal{L}(\mathrm{E})\) 。设 \(U \subset C\) 为 \(\mathbf{C}-\mathrm{Sp}(u)\) 的无界连通分量的一个开子集。若 M \textgreater{} 0 使得 \(\|R(u, \lambda)\| \leqslant M\) 对所有 \(\lambda \in U\) 成立，则 \(\|R(u, \lambda)\| < M\) 对所有 \(\lambda \in U\) 成立。特别地，若 \(\mathbf{C}-\mathrm{Sp}(u)\) 连通，则 u 的预解式的范数不可能在 \(\mathbf{C}-\mathrm{Sp}(u)\) 的非空开子集上取常值。
\end{enumerate}

\begin{enumerate}
\def\labelenumi{\arabic{enumi})}
\setcounter{enumi}{17}
\tightlist
\item
  设 E 为由复数组成的有界序列 \((x_{k})_{k\in\mathbf{Z}}\) 构成的复向量空间。对 E 中每个 \(x=(x_{k})\) ，我们记
\end{enumerate}

\[
\| x \| _ {\mathrm{E}} = | x _ {0} | + \sup _ {k \in \mathbf {Z} - \{0 \}} | x _ {k} |.
\]

\begin{enumerate}
\def\labelenumi{\alph{enumi})}
\item
  映射 \(x \mapsto \|x\|_{\mathrm{E}}\) 是 E 上的一个范数；赋以此范数的空间 E 是巴拿赫空间。
\item
  设 M \textgreater{} 2 为实数。设 \(\lambda = (\lambda_k)_{k \in \mathbf{Z}}\) 为由 \(\lambda_0 = M^{-1}\) 与 \(\lambda_k = 1\) （\(k \neq 0\) ）定义的序列。我们用 u 表示由 \(u((x_k)_{k \in \mathbf{Z}}) = (\lambda_k x_{k+1})_{k \in \mathbf{Z}}\) 定义的从 E 到 E 的线性映射。它是 E 的一个连续自同态。
\item
  以 0 为心、1/M 为半径的开圆盘包含于 u 的预解集。
\end{enumerate}

\(d)\) 对每个 \(\lambda\in\mathbf{C}\) （满足 \(|\lambda|<\inf(M^{-1},\frac{1}{2}-M^{-1})\) ），我们有 \(\|\mathrm{R}(u,\lambda)\|\leqslant\mathrm{M}\) 。（利用 IV 第 244 页公式 (6)。）

\begin{enumerate}
\def\labelenumi{\alph{enumi})}
\setcounter{enumi}{4}
\tightlist
\item
  设 \(e_{0} \in E\) 为序列 \((x_{k})_{k \in \mathbf{N}}\) ，它由 \(x_{0} = 1\) 与 \(x_{k} = 0\) （\(k \neq 0\) ）定义。设 \(\lambda \in C\) 满足 \(|\lambda| < \inf(M^{-1}, \frac{1}{2} - M^{-1})\) 。我们有 \(\|R(u, \lambda)e_{0}\|_{E} = M\) ，从而 \(\|R(u, \lambda)\| = M\) 。
\end{enumerate}

（本例与下一习题中的例子都归于 E. Shargorodsky，“On the level sets of the resolvent norm of a linear operator”，Bull. London Math. Soc. 40 (2008)，493–504。）

\begin{enumerate}
\def\labelenumi{\arabic{enumi})}
\setcounter{enumi}{18}
\tightlist
\item
  设 E 为希尔伯特空间 \(\ell_2(\mathbf{N}^*)\) 。设 \((\alpha_n)_{n \in \mathbf{N}^*}\) 为一列 \(\geqslant 2\) 的实数，使得 \(\alpha_n\) 趋于 \(+\infty\) （当 \(n \to +\infty\) ）。令 \(\beta_n = 1 + \alpha_n^{-1}\) （对所有 \(n \in \mathbf{N}^*\) ）。
\end{enumerate}

\begin{enumerate}
\def\labelenumi{\alph{enumi})}
\tightlist
\item
  设 \(\mathrm{F}\subset \mathrm{E}\) 为由满足下列条件的序列 \((x_{k})_{k\in \mathbf{N}^{*}}\in \mathrm{E}\) 组成的空间：级数
\end{enumerate}

\[
\sum_ {n \in \mathbf {N} ^ {*}} \alpha_ {n} ^ {2} | x _ {2 n} | ^ {2}
\]

收敛。它是 E 的一个稠密子空间。

\begin{enumerate}
\def\labelenumi{\alph{enumi})}
\setcounter{enumi}{1}
\tightlist
\item
  存在 E 上以 F 为定义域的一个闭部分算子 \(u\) ，使得对 F 中每个 \(x = (x_{k})_{k \in \mathbf{N}^{*}}\) ，我们有 \(u(x) = (y_{k})\) ，其中
\end{enumerate}

\[
\binom{y _ {2 k - 1}}{y _ {2 k}} = \left( \begin{array}{c c} 0 & \alpha_ {k} \\ \beta_ {k} & 0 \end{array} \right) \binom{x _ {2 k - 1}}{x _ {2 k}}
\]

对所有 \(k \in \mathbf{N}^*\) 成立。

\begin{enumerate}
\def\labelenumi{\alph{enumi})}
\setcounter{enumi}{2}
\tightlist
\item
  设 \(u_{n}\) 为自同态
\end{enumerate}

\[
x \mapsto \left( \begin{array}{c c} 0 & \alpha_ {n} \\ \beta_ {n} & 0 \end{array} \right) x
\]

（配备典范内积的希尔伯特空间 \(C^{2}\) 上的）。对每个 \(\lambda \in C\) （满足 \(|\lambda| < 1\) ），\(u_{n}\) 在 \(\lambda\) 处的预解式有定义，并且我们有

\[
\lim _ {n \to + \infty} \| \mathrm{R} (u _ {n}, \lambda) \| = 1.
\]

\begin{enumerate}
\def\labelenumi{\alph{enumi})}
\setcounter{enumi}{3}
\tightlist
\item
  对每个 \(n \in N^{*}\) 与每个 \(\lambda \in C\) （满足 \(|\lambda| < \frac{1}{2}\) ），我们有 \(\|R(u_{n}, \lambda)\| < 1\) 。
\item
  以 0 为心、半径 \(\frac{1}{2}\) 的开圆盘包含于 u 的预解集，且我们有 \(\|R(u,\lambda)\|=1\) 对所有 \(\lambda\in C\) （满足 \(|\lambda|<\frac{1}{2}\) ）成立。
\end{enumerate}

\begin{enumerate}
\def\labelenumi{\arabic{enumi})}
\setcounter{enumi}{19}
\tightlist
\item
  若 \(a\) 与 \(b\) 为满足 \(a < b\) 的实数，我们用 \(\mathcal{C}^2([a,b])\) 表示 \([a,b]\) 上连续、且其限制到 \(]a,b[\) 上为 \(\mathrm{C}^2\) 类的函数所成的向量空间。我们用 \(\mu\) 表示勒贝格测度。
\end{enumerate}

\begin{enumerate}
\def\labelenumi{\alph{enumi})}
\tightlist
\item
  设 \(q_{1}\) 与 \(q_{2}\) 为 \(\mathcal{C}^{2}([a,b])\) 中满足 \(q_{1} \geqslant q_{2}\) 的实值函数。对 \(i \in \{1,2\}\) ，设 \(f_{i} \in \mathcal{C}^{2}([a,b])\) 为方程
\end{enumerate}

\[
f _ {i} ^ {\prime \prime} + q _ {i} f _ {i} = 0.
\]

假设 \(f_{2}(a) = f_{2}(b) = 0\) 且 \(f_{2} \neq 0\) 。若函数 \(f_{1}\) 在 \(]a, b[\) 中不取零值，则 \(q_{1} = q_{2}\) ，且存在 \(c \in \mathbf{C}\) 使 \(f_{1} = cf_{2}\) 。（化归到 \(f_1\) 与 \(f_2\) 在 \(]a, b[;\) 上严格为正的情形；考虑朗斯基行列式 \(w = f_1f_2' - f_1'f_2\) ，先证明 \(w(a) = w(b)\) ，再证明 \(w = 0.\) ）

\begin{enumerate}
\def\labelenumi{\alph{enumi})}
\setcounter{enumi}{1}
\tightlist
\item
  设 \(q \in \mathcal{C}^2([0,1])\) 满足 \(q \geqslant 0\) 。我们用 \(\widetilde{\mathrm{SL}}_q\) 表示 \(\mathrm{L}^2([0,1],\mu)\) 上以
\end{enumerate}

\[
\mathrm{D} = \{f \in \mathcal {C} ^ {2} ([ 0, 1 ]) \mid f (0) = f (1) = 0 \}
\]

为定义域、满足

\[
\widetilde {\mathrm{SL}} _ {q} (f) = - f ^ {\prime \prime} + q f
\]

的部分微分算子（斯图姆–刘维尔算子）。算子 \(\widetilde{SL}_{q}\) 是对称的且本质自伴。我们用 \(SL_{q}\) 表示 \(\widetilde{SL}_{q}\) 的闭包。

\begin{enumerate}
\def\labelenumi{\alph{enumi})}
\setcounter{enumi}{2}
\tightlist
\item
  设 \(\lambda \in \mathrm{Sp}(\mathrm{SL}_q)\) 为 \(\mathrm{SL}_q\) 的一个特征值。与 \(\lambda\) 相伴的特征空间维数为 1，且由一个属于 D 的实值函数生成。
\end{enumerate}

\(d\) ) \(\lambda\) 作为 \(\mathrm{SL}_q\) 的特征值，满足 \(\lambda \geqslant \pi^2\) 。（先证明 \(\lambda\) 是正的：利用问题 \(a\) ）取 \(q_{1} = q\) 与 \(q_{2} = 0\) 。）

\begin{enumerate}
\def\labelenumi{\alph{enumi})}
\setcounter{enumi}{4}
\tightlist
\item
  对 s \textgreater{} 0，存在唯一的函数 \(f_{s} \in \mathcal{C}^{2}([0,1])\) 是线性微分方程
\end{enumerate}

\[
- f _ {s} ^ {\prime \prime} + q f - s ^ {2} f = 0
\]

的解，且满足 \(f_{s}(0)=0\) 与 \(f_{s}^{\prime}(0)=s\) 。令 \(f(x,s)=f_{s}(x)\) ；映射 f 是 \(C^{2}\) 类函数，定义在 {]}0,1{[}× \(R_{+}^{*}\) 上。\(SL_{q}\) 的特征值形如 \(\lambda=s^{2}\) ，其中 s 满足 \(f(1,s)=0\) 。

\(f)\) 对 \(t \in [0,1]\) 与 \(s > 0\) ，设 \(\mathrm{N}(t,s) \in \overline{\mathbf{N}}\) 为 \(x \in [0,t]\) 中满足 \(f(x,s) = 0\) 者的个数。我们有 \(\mathrm{N}(t,s) \in \mathbf{N}\) ，且函数 \(s \mapsto \mathrm{N}(t,s)\) 对每个固定的 \(t\) 是递增的。（利用问题 \(a\) ）。）

\(g)\) 设 \(t \in [0,1]\) 且 \(s_0 > 0\) 。设 \(m = \mathrm{N}(t,s_0) \in \mathbf{N}\) 。用

\[
0 = x _ {1} (s _ {0}) <   x _ {2} (s _ {0}) <   \dots <   x _ {m} (s _ {0})
\]

表示 \(f_{s_0}\) 在 \([0,t]\) 中的零点。对 \(s \geqslant 0\) ，用

\[
0 = x _ {1} (s) <   \dots <   x _ {m} (s)
\]

表示前 \(m\) 个零点（\(f_{s}\) 在 \([0,t]\) 中的）。函数 \(s\mapsto x_{j}(s)\) 是 \(\mathbf{C}^1\) 类的，并且我们有

\[
\partial_ {1} (x _ {j} (s), s) x _ {j} ^ {\prime} (s) = - \partial_ {2} f (x _ {j} (s), s).
\]

\begin{enumerate}
\def\labelenumi{\alph{enumi})}
\setcounter{enumi}{7}
\tightlist
\item
  \(SL_{q}\) 的特征值之集是某个严格正实数递增序列的取值集
\end{enumerate}

\[
0 <   \lambda_ {0} <   \lambda_ {1} <   \dots <   \lambda_ {n} <   \dots
\]

该序列趋于 \(+\infty\) 。此外，与 \(\lambda_{n}\) 相应的特征空间维数为 1；若 \(\varphi_{n}\) 是它的一个基，则 \(\varphi_{n} \in \mathcal{C}^{2}([0,1])\) ，且函数 \(\varphi_{n}\) 恰在 \(n\) 个不同点处（\([0,1]\) 中）取零。（为证明 \(\mathrm{SL}_{q}\) 存在无穷多个特征值，利用问题 \(a\) ）取 \(q_{1} = M - s^{2}\) ，其中 \(M \in R\) 选得适当，验证 \(N(1, s)\) 趋于 \(+\infty\) （当 \(s \to +\infty\) ）。）

\begin{enumerate}
\def\labelenumi{\arabic{enumi})}
\setcounter{enumi}{20}
\tightlist
\item
  我们沿用上一习题的记号。我们用 \(\varphi_{n}\) 表示 \(\mathrm{SL}_{q}\) 的属于特征值 \(\lambda_{n}\) 且满足 \(\varphi_{n}^{\prime}(0)=\lambda_{n}\) 的唯一特征函数。
\end{enumerate}

\begin{enumerate}
\def\labelenumi{\alph{enumi})}
\tightlist
\item
  设 \(\mathrm{M} = \sup \{q(x)\}\) 。对每个 \(n \in \mathbf{N}\) ，我们有
\end{enumerate}

\[
\pi^ {2} (n + 1) ^ {2} \leqslant \lambda_ {n} \leqslant \pi^ {2} (n + 1) ^ {2} + \mathrm{M},
\]

并且此外

\[
\lambda_ {n} = \pi^ {2} (n + 1) ^ {2} \Bigl (1 + \mathrm{O} \Bigl (\frac {1}{n ^ {2}} \Bigr) \Bigr)
\]

当 \(n \rightarrow +\infty\) 时成立。

\begin{enumerate}
\def\labelenumi{\alph{enumi})}
\setcounter{enumi}{1}
\tightlist
\item
  设 \(s > 0\) 。设 \(f \in \mathcal{C}^2([0,1])\) 满足
\end{enumerate}

\[
- f ^ {\prime \prime} + s ^ {2} f = - q f, \qquad f (0) = 0, \quad f ^ {\prime} (0) = s.
\]

于是我们有 \((1-s^{-1}u_{s})f=g_{s}\) ，其中 \(g_{s}(x)=\sin(sx)\) ，而 \(u_{s}\) 是 \(\mathcal{C}([0,1])\) 上由下式给出的积分算子：

\[
u _ {s} (f) (x) = \int_ {0} ^ {x} \sin (s (x - y)) q (y) f (y) d \mu (y).
\]

\begin{enumerate}
\def\labelenumi{\alph{enumi})}
\setcounter{enumi}{2}
\item
  当 \(s\) 充分大时，我们有 \(s \notin \operatorname{Sp}(u_s)\) 。
\item
  对每个固定的 x，我们有
\end{enumerate}

\[
\varphi_ {n} (x) = \sin ((n + 1) x) + \mathrm{O} (n ^ {- 1})
\]

当 \(n \rightarrow +\infty\) 时成立。

\begin{enumerate}
\def\labelenumi{\arabic{enumi})}
\setcounter{enumi}{21}
\tightlist
\item
  我们沿用上一习题的记号。
\end{enumerate}

\begin{enumerate}
\def\labelenumi{\alph{enumi})}
\tightlist
\item
  存在函数 \(f_{1}\) 与 \(f_{2}\) 属于 \(\mathcal{C}^{2}([a,b])\) ，满足 \(-f_{i}^{\prime\prime} + qf_{i} = 0\) 且
\end{enumerate}

\[
\begin{array}{c c} f _ {1} (0) = 0, & f _ {1} ^ {\prime} (0) = 1 \\ f _ {2} (1) = 0, & f _ {2} ^ {\prime} (1) = - 1. \end{array}
\]

这些函数是这些方程的唯一解，且函数 \(w = f_{1}f_{2}' - f_{1}'f_{2}\) 是常数。

\begin{enumerate}
\def\labelenumi{\alph{enumi})}
\setcounter{enumi}{1}
\tightlist
\item
  设 \(g \in \mathcal{C}([0,1])\) 。存在唯一的函数 \(\varphi \in \mathcal{C}^2([0,1])\) 满足
\end{enumerate}

\[
- \varphi^ {\prime \prime} + q \varphi = g,
\]

并且我们有 \(\varphi = c_{1}f_{1} + c_{2}f_{2}\) ，其中

\[
c _ {1} (x) = - \frac {1}{w} \int_ {x} ^ {1} g (y) f _ {2} (y) d \mu (y), \qquad c _ {2} (x) = - \frac {1}{w} \int_ {0} ^ {x} g (y) f _ {1} (y) d \mu (y)
\]

对所有 \(x \in [0,1]\) 成立。

\begin{enumerate}
\def\labelenumi{\alph{enumi})}
\setcounter{enumi}{2}
\tightlist
\item
  部分算子 \(\mathrm{SL}_q\) 是满射的；其逆是映射 \(g\mapsto \varphi\) ，且存在函数 \(\mathrm{G}\in \mathcal{C}([0,1]^2)\) 使得
\end{enumerate}

\[
\mathrm{SL} _ {q} ^ {- 1} (g) (x) = \int_ {0} ^ {1} \mathrm{G} (x, y) g (y) d \mu (y)
\]

对所有 \(g \in \mathcal{C}([0,1])\) 与几乎所有 \(x \in [0,1]\) 成立。我们有 \(G(x,y) = G(y,x)\) 对所有 \((x,y) \in [0,1]^2\) 成立。

\(d)\) 算子 \(\mathrm{SL}_q^{-1}\) 在 \(\mathrm{L}^2 ([0,1],\mu)\) 上是紧的且单射的。

\begin{enumerate}
\def\labelenumi{\alph{enumi})}
\setcounter{enumi}{4}
\tightlist
\item
  诸函数 \((\varphi_{n})\) 构成 \(\mathrm{L}^{2}([0,1])\) 的一个标准正交基，且 \(SL_{q}\) 的谱与其特征值之集重合。
\end{enumerate}

\begin{enumerate}
\def\labelenumi{\arabic{enumi})}
\setcounter{enumi}{22}
\tightlist
\item
  设 \(E = \ell_{\mathbf{C}}^{2}(\mathbf{N})\) ，u 为由
\end{enumerate}

\[
u ((x _ {i}) _ {i \in \mathbf {N}}) = (0, x _ {1}, \dots , x _ {n}, \dots)
\]

对所有 \((x_{i})\in \mathrm{E}\) 定义的自同态。

\begin{enumerate}
\def\labelenumi{\alph{enumi})}
\item
  \(u\) 的谱是 \(\mathbf{C}\) 中的闭单位圆盘。
\item
  u 的剩余谱是 C 中的开单位圆盘。
\item
  \(u\) 的本质谱是 \(\mathbf{C}\) 中的单位圆周。
\end{enumerate}

\begin{enumerate}
\def\labelenumi{\arabic{enumi})}
\setcounter{enumi}{23}
\tightlist
\item
  设 E 为复希尔伯特空间，u 为 E 上的一个对称部分算子。设 \(F = E \oplus E\) ，v 为 F 上以 \(\text{dom}(u^{*}) \oplus \text{dom}(\overline{u})\) 为定义域、满足 \(v(x, y) = (\overline{u}(y), u^{*}(x))\) 的部分算子。
\end{enumerate}

\begin{enumerate}
\def\labelenumi{\alph{enumi})}
\item
  部分算子 v 是自伴的。
\item
  设 \(\widetilde{\mathrm{E}}\) 为线性映射 \(\delta\colon x\mapsto(x,x)\) 的像。\(v\) 到 \(\widetilde{\mathrm{E}}\) 的约化是以 \(\delta(\mathrm{dom}(u))\) 为定义域、由 \((x,x)\mapsto(u(x),u(x))\) 给出的部分算子。
\end{enumerate}

